\section{可操作核查清单（C-4）}\label{sec:checklist}

\subsection{设计原则}

清单的目的是把前两节的框架落到可执行的检查动作上。设计遵循三条原则：\textbf{每项可判定}——给定一条推理链，检查者能给出通过/不通过；\textbf{每项对应一类失效}——不做无指向的“再检查一遍”；\textbf{每项代价已知}——检查者能判断何时值得做。

\begin{table}[t]
\centering
\small
\caption{十项核查清单。末列为该检查主要针对的失效模式或指标。}
\label{tab:checklist}
\begin{tabularx}{\textwidth}{@{}c p{3.3cm} X c@{}}
\toprule
\textbf{\#} & \textbf{检查项} & \textbf{判定问题} & \textbf{针对} \\
\midrule
1 & 题面歧义枚举 & 题面是否存在两种以上合理读法？若存在，选中的解释是否显式记录？ & F1 \\
2 & 假设集完备性 & $\sigma$ 的 $H$ 是否覆盖了题面全部显式条件与选定解释的隐含前提？ & F1, F4 \\
3 & 变量类型标注 & 每个变量是否带类型与定义域？有无未标注的自由变量？ & F2 \\
4 & 定义域边界 & 结论成立的范围是否与变量定义域一致？是否在边界点单独检查？ & F2 \\
5 & 量词范围 & 全称/存在量词的作用域是否与自然语言原句一致？ & F2 \\
6 & 反译可读性 & 能否由 $\sigma$ 逐句还原出与原推理链语义一致的自然语言陈述？ & F3 \\
7 & 步骤前提核对 & 每一步 $s_i$ 所依赖的假设是否都在 $H$ 中，且未被后续步骤静默加强？ & F4 \\
8 & 反例尝试 & 是否主动构造过至少一个试图推翻结论的实例（含边界值、退化情形）？ & $\mathrm{CDR}$ \\
9 & 验证层级声明 & 本次结论达到的验证层级是否被明确标注？是否与结论强度匹配？ & $V$-rate \\
10 & 未通过项留痕 & 未通过的检查是否被记录并作为结论的适用条件，而非直接删除？ & 全部 \\
\bottomrule
\end{tabularx}
\end{table}

\subsection{使用方式与优先级}

清单按三档使用。\textbf{轻档（第 1、3、6 项）}适用于任何一次模型给出的数学推理，代价最低，能拦住 F1/F2/F3 中最常见的表现。\textbf{常规档（第 1--7、9 项）}适用于结论会被下游使用、但无需形式化内核的场景：此时应显式声明验证层级，并把未通过项写进结论的适用条件。\textbf{完整档（全部十项）}适用于结论将进入形式化验证或作为后续推导依据的场景，其中第 8 项（反例尝试）与第 6 项（反译可读性）是能否归因的关键。

第 10 项值得单独说明。实践中一个反复出现的反模式是：检查发现某一步的前提不成立后，直接在正文中删掉该步或改写结论，而不保留“该结论仅在额外假设下成立”这一信息。这样做的结果是后续使用者拿到一个无条件的结论，其成立条件已经不可追。留痕的成本极低，但它把“当前无法证明”与“已被证伪”区分开，避免了不可逆的信息损失。

\subsection{与其他路线的关系}

本清单不替代形式化验证，也不要求使用者具备证明助手经验：第 1--7、9、10 项都可以在自然语言与轻量符号表示之间完成。它与形式化路线的关系是分工——清单负责在进入内核之前把 F1--F4 尽量拦下，从而降低“形式化一个错误问题”的概率；内核负责在清单通过之后提供机器可检的判定。反过来，清单也不替代答案级检查：第 5 项与第 4 项的正确执行通常需要符号计算或数值回代来确认（见第 \ref{sec:case} 节例 2）。

\section{案例走查}\label{sec:case}

本节用两个可手工核对的小规模例子说明框架如何工作。两例均为为说明机制而构造的教学例子，不引用任何外部数据集，也不作经验性统计主张。

\subsection{例 1：歧义未消解（F1）}

\textbf{题面。}“三角形 $ABC$ 的面积为 $12$，$AB = 6$，求 $AB$ 边上的高。”

\textbf{L1 输出（节选）。}由面积公式 $S = \frac{1}{2}\cdot \text{底} \cdot \text{高}$ 得 $12 = \frac{1}{2}\cdot 6 \cdot h$，解得 $h = 4$。

\textbf{失效定位。}此处“$AB$ 边上的高”在题面中只有一种自然读法，本例的歧义在别处：面积公式中的“高”与题面中的“$AB$ 边上的高”是否同一个量，取决于 $AB$ 是否被理解为底边。若题面写成“三角形 $ABC$ 的面积为 $12$，$BC = 6$，求 $A$ 到 $BC$ 的距离”，则不存在歧义。当题面代词指代不明（“该边”“此高”）时，规约层会静默选定一个所指。

\textbf{按清单走查。}第 1 项（题面歧义枚举）：检查题面中的每处指代是否唯一确定 → 若存在两种读法，第 2 项要求把选中的读法写入 $H$。若 $\sigma$ 的 $H$ 中只有 $S = 12$、$|AB| = 6$，而不含“$AB$ 为所求高对应的底边”，则第 2 项不通过，结论 $h = 4$ 应被标注为“在选定解释下成立”。

\textbf{失效的可检出性。}此例中 V0（答案级）会核对最终数值，若参考答案同为 $4$ 则整条链条被视为通过，F1 完全不可见。这正是 F1 之所以危险的原因：它的后果不是数值错误，而是结论的\emph{适用条件}被丢失。

\subsection{例 2：类型不完整与语义漂移（F2、F4）}

\textbf{题面。}“求所有正整数 $n$ 使得 $n^2 - n$ 能被 $6$ 整除。”

\textbf{L1 输出（节选）。}$n^2 - n = n(n-1)$。连续两整数之积必为偶数；又 $n(n-1)(n-2)$ 为三连续整数之积故被 $3$ 整除。但此时我们只需 $3 \mid n(n-1)$：当 $n \equiv 0,1 \pmod 3$ 时成立，当 $n \equiv 2 \pmod 3$ 时 $n-1 \equiv 1$、$n \equiv 2$，均不被 $3$ 整除。故结论为 $n \equiv 0$ 或 $1 \pmod 3$。

\textbf{失效定位。}该推导的最终结论是正确的，但中间出现两次问题。其一（F2）：$n(n-1)(n-2)$ 被 $3$ 整除不能推出 $n(n-1)$ 被 $3$ 整除——这一步的变量类型信息（这里的对象是整数环上的因子，而非任意实数）被忽略，作者随后补上的修正恰好说明这一点。其二（F4）：修正过程用“当 $n\equiv 2 \pmod 3$ 时 $n-1\equiv 1$、$n\equiv 2$”来排除，其中隐含使用了 $n$ 的正性以外的条件；若把 $n$ 的范围改到 $\mathbb{Z}$，结论的形式不变但论证需重写。

\textbf{按清单走查。}第 3 项（变量类型标注）：$n$ 的类型为 $\mathbb{Z}^{+}$，则所有模运算的合法性需在 $\mathbb{Z}$ 上成立 → 通过，但不等于结论在 $\mathbb{Z}$ 上成立，这一点应写入第 4 项（定义域边界）：需单独检查 $n \le 0$ 的情形。第 7 项（步骤前提核对）：检查“$3 \mid n(n-1)$”这一结论所依赖的前提是否在 $H$ 中 → 该前提并未由题面给出，也未由前几步推出，故不通过。第 8 项（反例尝试）：取 $n = 2$，得 $n^2 - n = 2$，不被 $6$ 整除 → 与最终结论一致；取 $n = 3$，得 $6$，被 $6$ 整除 → 一致；取 $n = 4$，得 $12$，被 $6$ 整除 → 一致。这说明结论本身正确，但中间推理存在断裂——即“结论对、推理不成立”，这类情形必须由步骤级或结构化级检查捕获，答案级比对无能为力。

\textbf{验证层级的差异。}若把此题送入 V3：需先把 $\sigma$ 编译为“$\forall n \in \mathbb{Z}^{+},\ 6 \mid n^2 - n \iff n \bmod 3 \in \{0,1\}$”并交内核判定。此时 $\sigma$ 中若已丢失 $n$ 的正性（F2）或错误地保留了断裂的中间步骤（F4），内核会忠实地验证一个与题面不同的命题，并返回 pass。这说明 R1：\emph{内核的强度不能补偿规约的缺陷}。

\subsection{两例的共同结论}

两个例子分别落在“结论对但适用条件丢失”（例 1）与“结论对但推理断裂”（例 2）两种情形上。二者都无法被答案级验证检出，都必须依赖对 $\sigma$ 的结构化检查（清单第 2--7 项）才能暴露；而第 8 项的反例尝试虽不能定位断裂位置，却是区分“验证流程有效”与“验证流程走过场”的最低成本手段。
